\documentclass[10pt,a4paper]{article}
\usepackage[margin=20mm,headheight=15pt,headsep=7mm,footskip=10mm]{geometry}
\usepackage{fontspec,amsmath,amssymb,mathtools,microtype}
\setmainfont{lmroman10-regular.otf}[BoldFont=lmroman10-bold.otf,ItalicFont=lmroman10-italic.otf,BoldItalicFont=lmroman10-bolditalic.otf]
\setsansfont{lmsans10-regular.otf}[BoldFont=lmsans10-bold.otf,ItalicFont=lmsans10-oblique.otf,BoldItalicFont=lmsans10-boldoblique.otf]
\setmonofont[Scale=.86]{lmmono10-regular.otf}[ItalicFont=lmmono10-italic.otf]
\usepackage{ragged2e}\usepackage{xcolor,booktabs,longtable,array,graphicx,enumitem,fancyhdr,needspace,titlesec,hyperref,xurl}
\definecolor{AtlasNavy}{HTML}{18334C}\definecolor{AtlasTeal}{HTML}{116D73}\definecolor{AtlasGray}{HTML}{52616D}
\hypersetup{colorlinks=true,linkcolor=AtlasTeal,urlcolor=AtlasTeal,pdfcreator={Complexity Atlas},unicode=true}
\urlstyle{same}\setlength{\parindent}{0pt}\setlength{\parskip}{4.8pt}\setlength{\emergencystretch}{3em}
\renewcommand{\baselinestretch}{1.08}\setcounter{secnumdepth}{0}\setlist{nosep,leftmargin=1.5em}
\titleformat{\section}{\Large\bfseries\sffamily\color{AtlasNavy}\raggedright}{}{0em}{}
\titleformat{\subsection}{\large\bfseries\sffamily\color{AtlasTeal}\raggedright}{}{0em}{}
\titleformat{\subsubsection}{\normalsize\bfseries\color{AtlasNavy}\raggedright}{}{0em}{}
\titlespacing*{\section}{0pt}{14pt}{6pt}\titlespacing*{\subsection}{0pt}{10pt}{5pt}
\pagestyle{fancy}\fancyhf{}\fancyhead[L]{\small\sffamily\color{AtlasGray}COMPLEXITY ATLAS}
\fancyhead[R]{\small\sffamily\color{AtlasGray}ENGLISH EDITION}\fancyfoot[C]{\small\thepage}\renewcommand{\headrulewidth}{.3pt}
\newcommand{\AtlasMath}[1]{\par\begingroup\centering\setbox0=\hbox{$\displaystyle #1$}\ifdim\wd0>1.4\linewidth\typeout{ATLAS WIDE FORMULA: \the\wd0}\fi\ifdim\wd0>\linewidth\resizebox{\linewidth}{!}{\box0}\else\box0\fi\par\endgroup}
\newcommand{\AtlasInlineMath}[1]{\begingroup\setbox0=\hbox{$#1$}\ifdim\wd0>.4\linewidth\endgroup\AtlasMath{#1}\else\endgroup$#1$\fi}
\pdfstringdefDisableCommands{\def\AtlasInlineMath#1{#1}}
\DeclareMathOperator{\polylog}{polylog}\DeclareMathOperator{\dist}{dist}\DeclareMathOperator{\diam}{diam}
\DeclareMathOperator{\rank}{rank}\DeclareMathOperator{\tr}{tr}\DeclareMathOperator{\diag}{diag}
\DeclareMathOperator{\prox}{prox}\DeclareMathOperator{\KL}{KL}\DeclareMathOperator{\Var}{Var}
\providecommand{\eps}{\varepsilon}\providecommand{\R}{\mathbb R}\providecommand{\E}{\mathbb E}
\providecommand{\mathbbm}[1]{\mathbf{#1}}\providecommand{\mathds}[1]{\mathbf{#1}}
\providecommand{\llbracket}{[\![}\providecommand{\rrbracket}{]\!]}
\begin{document}
\raggedright

{\LARGE\bfseries\sffamily\color{AtlasNavy} Unknown-tail stochastic nonconvex stationarity with known smoothness and initial gap\par}\vspace{5pt}


\textbf{Setting ID.} \texttt{SO\_\allowbreak{}FO\_\allowbreak{}S03\_\allowbreak{}unknown\_\allowbreak{}alpha\_\allowbreak{}known\_\allowbreak{}scales}



Single-objective optimization



A common algorithm adapts to an unknown moment exponent and noise scale when population smoothness and initial-gap budgets are known.



\textbf{Status.} Logarithmic gap; Bounds missing



\textbf{TASK tags.} solution concept: Gradient stationarity; task family: Single objective



\textbf{MODEL tags.} smoothness: Lipschitz gradient



\textbf{ORACLE tags.} feedback quality: Stochastic unbiased conditional moment; operation: Stochastic gradient; public information: Hidden objects explicitly declared; public information: Public inputs explicitly declared; scalar derivative order: Scalar-function first derivative



\textbf{ALGORITHMS tags.} randomization: Private randomness allowed; restriction: No additional span restriction explicitly imposed



\textbf{GUARANTEE tags.} mode: Expectation; quantifier: Algorithm before instance



\textbf{RESOURCE tags.} budget: Fixed worst-case budget; measure: N gradient sample



\needspace{7\baselineskip}\section{1. Setting specification}


\needspace{4\baselineskip}\subsection{Task}


Construct one budget-indexed algorithm family for smooth nonconvex stationarity that adapts to an unknown centered-noise moment exponent.



\needspace{4\baselineskip}\subsection{Admissible model}


The differentiable lower-bounded \AtlasInlineMath{F:\mathbb R^d\to\mathbb R} satisfies \AtlasInlineMath{F(x_0)-\inf F\le\Delta} and \AtlasInlineMath{\operatorname{Lip}(\nabla F)\le L}. The positive budgets \AtlasInlineMath{L,\Delta}, initial point and finite dimension are public. For some fixed \AtlasInlineMath{p\in(1,2]} and \AtlasInlineMath{\sigma\ge0}, the oracle noise has centered conditional pth moment at most \AtlasInlineMath{\sigma^p}. Neither p nor sigma is supplied to the algorithm.



\needspace{4\baselineskip}\subsection{Oracle interface}


A query at a history-measurable point returns only a full vector \AtlasInlineMath{g_t} with \AtlasInlineMath{\mathbb E[g_t\mid\mathcal F_{t-1}]=\nabla F(x_t)} and \AtlasInlineMath{\mathbb E[\|g_t-\nabla F(x_t)\|_2^p\mid\mathcal F_{t-1}]\le\sigma^p}. The conditional law may depend on the history. No function value, repeated evaluation of one sample at different points, or sample-gradient smoothness is available.



\needspace{4\baselineskip}\subsection{Output}


Return an arbitrary measurable point \AtlasInlineMath{\widehat x_T} after at most T charged gradient observations.



\needspace{4\baselineskip}\subsection{Loss or accuracy criterion}


The risk is \AtlasInlineMath{\mathcal R_T^{\mathsf A}(p;L,\Delta,\sigma)=\sup_{d,I,\mathcal O}\mathbb E\|\nabla F(\widehat x_T)\|_2}, with one common family \AtlasInlineMath{\mathsf A=(\mathsf A_T)} required to apply for every fixed p. Constants in the analyzed risk may depend on p; the algorithm does not.



\needspace{4\baselineskip}\subsection{Resource accounting}


T is a public deterministic cap on every path. Training, candidate generation, validation and selection calls all count. Pointwise inversion of a risk bound at a specified p is an analysis quantity and does not furnish a stopping rule that knows p.



\needspace{4\baselineskip}\subsection{Algorithm views and success guarantees}


Queries, stopping decisions and outputs are arbitrary measurable functions of the observed history and public inputs. No linear-span or zero-respecting restriction is imposed. Internal exact-real computation and storage are uncharged.



A private random seed independent of the fixed instance is permitted.



The complete algorithm family is the same for all p and sigma. Known smoothness and gap budgets may enter its step sizes; unknown tail and noise scales may not.




\AtlasMath{\mathcal R_T^{\mathsf A}(p;L,\Delta,\sigma)=\sup_{d,I,\mathcal O}\mathbb E\|\nabla F(\widehat x_T)\|_2}




Choose one common family with inputs only T,L,Delta,x0 and d. For every fixed p, evaluate its risk uniformly over the corresponding instance and oracle class. The family may not encode a p-dependent tuning choice.



\textbf{Budget convention.} fixed\_worst\_case\_cap



\needspace{7\baselineskip}\section{2. Notation and parameters}


\AtlasInlineMath{d}: Public finite ambient dimension \textbf{Domain:} \AtlasInlineMath{d\in\mathbb N_{\ge1}}.



\AtlasInlineMath{T}: Public total sample cap \textbf{Domain:} \AtlasInlineMath{T\in\mathbb N_{\ge1}}.



\AtlasInlineMath{L}: Known population-gradient smoothness bound \textbf{Domain:} \AtlasInlineMath{L>0}.



\AtlasInlineMath{\Delta}: Known initial gap bound \textbf{Domain:} \AtlasInlineMath{\Delta>0}.



\AtlasInlineMath{p}: Unknown fixed centered-noise moment exponent \textbf{Domain:} \AtlasInlineMath{1<p\le2}.



\AtlasInlineMath{\sigma}: Unknown centered-noise scale \textbf{Domain:} \AtlasInlineMath{\sigma\ge0}.



\AtlasInlineMath{\varepsilon}: Tolerance used for pointwise risk inversion \textbf{Domain:} \AtlasInlineMath{\varepsilon>0}.



\AtlasInlineMath{a(p)}: Normalized risk exponent \textbf{Domain:} \AtlasInlineMath{a(p)>0}.



\AtlasInlineMath{a(p)=(p-1)/(3p-2)}



\AtlasInlineMath{b(p)}: Mean-estimation exponent \textbf{Domain:} \AtlasInlineMath{b(p)>0}.



\AtlasInlineMath{b(p)=(p-1)/p}



\AtlasInlineMath{n}: Training budget per candidate in the portfolio inequality \textbf{Domain:} \AtlasInlineMath{n\in\mathbb N,\ n\ge8}.



\AtlasInlineMath{m}: Fresh validation samples per candidate \textbf{Domain:} \AtlasInlineMath{m\in\mathbb N_{\ge1}}.



\AtlasInlineMath{M}: Number of dyadic candidates \textbf{Domain:} \AtlasInlineMath{M\in\mathbb N}.



\AtlasInlineMath{M=|\mathcal T_n|,\quad\mathcal T_n=\{2^j:j\in\mathbb Z_{\ge0},2^j\le n/2\}}



\needspace{12\baselineskip}\section{3. Complexity bracket}


\needspace{4\baselineskip}\subsection{Compared quantities}


\needspace{4\baselineskip}\subsubsection{Q1 - minimax metric resource quantity}


\AtlasInlineMath{\mathcal R_T^{\mathsf A}(p;1,1,1)}, with one common family for all p



\textbf{Scope.} The union of the explicit parameter and dimension domains of the linked comparison regimes; each bound retains its own applicability conditions.



\textbf{Optimization direction.} minimize



\textbf{Resource contract.} expression: \AtlasInlineMath{N_{\rm sample}\le T}; description: All candidate training and validation responses count toward the fixed cap.



\needspace{4\baselineskip}\subsubsection{Q2 - minimax metric resource quantity}


\AtlasMath{\mathcal R_T^{\mathsf A}(p;L,\Delta,\sigma)}


\textbf{Scope.} The union of the explicit parameter and dimension domains of the linked comparison regimes; each bound retains its own applicability conditions.



\textbf{Optimization direction.} minimize



\textbf{Resource contract.} expression: \AtlasInlineMath{M(n+m)\le T}; description: Separate per-candidate training and validation allocations are charged in full.



\needspace{4\baselineskip}\subsection{R1 - Normalized unknown-tail risk with known smoothness and gap}


\textbf{Region.} \AtlasInlineMath{L=\Delta=\sigma=1} for evaluation of the instance class, fixed \AtlasInlineMath{1<p\le2}, and all sufficiently large integer T (the sufficient threshold may depend on p). The algorithm is not given p or sigma. Upper bounds hold in every finite dimension. The lower bound permits a dimension depending on T; Corollary sample-complexity imports Liu, extended-version Theorem 11. No fixed-dimensional lower rate is asserted.



\textbf{Quantity.} Q1



\textbf{Resource.} All candidate training and validation responses count toward the fixed cap.



\textbf{Status.} Logarithmic gap



\textbf{Lower bound.} B2



\AtlasMath{\Omega_p\!\left(T^{-(p-1)/(3p-2)}\right)}


\textbf{Upper bound.} B1



\AtlasMath{O_p\!\left((\log(eT)/T)^{(p-1)/(3p-2)}\right)}


\textbf{Gap or comparison.} The risk ratio is at most \AtlasInlineMath{O_p((\log(eT))^{(p-1)/(3p-2)})}. Pointwise sample complexity is between \AtlasInlineMath{\Omega_p(\varepsilon^{-(3p-2)/(p-1)})} and \AtlasInlineMath{O_p(\varepsilon^{-(3p-2)/(p-1)}\log(e/\varepsilon))}.



\needspace{4\baselineskip}\subsection{R2 - General-scale portfolio oracle inequality}


\textbf{Region.} \AtlasInlineMath{L,\Delta>0}, \AtlasInlineMath{\sigma\ge0}, fixed \AtlasInlineMath{1<p\le2}, integers \AtlasInlineMath{n\ge8,m\ge1}, and \AtlasInlineMath{M(n+m)\le T}. The upper bound holds for every finite d.



\textbf{Quantity.} Q2



\textbf{Resource.} Separate per-candidate training and validation allocations are charged in full.



\textbf{Status.} Lower bound missing



\textbf{Lower bound.} No admitted bound for this comparison



\textbf{Upper bound.} B3



\AtlasMath{O_p\!\left(\min_{\tau\in\mathcal T_n}\{\sqrt{L\Delta\tau/n}+\sigma\tau^{-(p-1)/p}\}+\sigma M^{1/p}m^{-(p-1)/p}\right)}


\textbf{Gap or comparison.} Only the displayed upper bound is recorded for the full joint parameter domain.



\needspace{7\baselineskip}\section{4. Bounds and supporting evidence}


\needspace{4\baselineskip}\subsection{B1 - Dyadic minibatch normalized-gradient portfolio and fresh empirical-gradient validation}


\textbf{Admission.} Included in the bracket



\textbf{Direction.} upper



\textbf{Witness or lower-bound scope.} Randomized witness



\AtlasMath{O_p\!\left((\log(eT)/T)^{(p-1)/(3p-2)}\right)}


\textbf{Sources.} \href{https://github.com/WU-Wendao/Colombo/blob/main/colombo_sensorium/Complexity_altas/data/sources/colombo/N05/manuscript.pdf}{[1]}



\textbf{Conditions.} \AtlasInlineMath{L=\Delta=\sigma=1} for evaluation of the instance class, fixed \AtlasInlineMath{1<p\le2}, and all sufficiently large integer T (the sufficient threshold may depend on p). The algorithm is not given p or sigma.; The same family simultaneously satisfies the displayed pointwise risk curve for each fixed p.



\textbf{Dimension conditions.} Upper bounds hold in every finite dimension. The lower bound permits a dimension depending on T; Corollary sample-complexity imports Liu, extended-version Theorem 11. No fixed-dimensional lower rate is asserted.



\textbf{Constants.} Constants and sufficient large-budget thresholds may depend on fixed p; they are not uniform as p decreases to 1.



\textbf{Logarithmic factors.} A single logarithm in pointwise sample complexity; its displayed power in risk.



\textbf{Source location.} \href{https://github.com/WU-Wendao/Colombo/blob/main/colombo_sensorium/Complexity_altas/data/sources/colombo/N05/manuscript.pdf}{[1]} Theorem main-upper; equation main-upper; verified labels: sec:main-upper (line 544), thm:main-upper (line 569), eq:main-upper (line 572)



Dyadic minibatch normalized-gradient portfolio and fresh empirical-gradient validation. One randomized common algorithm family independent of p and sigma.



The same family simultaneously satisfies the displayed pointwise risk curve for each fixed p.



\needspace{4\baselineskip}\subsection{B2 - Known-index stochastic first-order lower bound, valid for the smaller class of tuning-free families}


\textbf{Admission.} Included in the bracket



\textbf{Direction.} lower



\textbf{Witness or lower-bound scope.} All randomized adaptive algorithms



\AtlasMath{\Omega_p\!\left(T^{-(p-1)/(3p-2)}\right)}


\textbf{Sources.} \href{https://github.com/WU-Wendao/Colombo/blob/main/colombo_sensorium/Complexity_altas/data/sources/colombo/N05/manuscript.pdf}{[1]}



\textbf{Conditions.} \AtlasInlineMath{L=\Delta=\sigma=1} for evaluation of the instance class, fixed \AtlasInlineMath{1<p\le2}, and all sufficiently large integer T (the sufficient threshold may depend on p). The algorithm is not given p or sigma.; The source lower construction obeys the centered moment convention and allows growing dimension.



\textbf{Dimension conditions.} Upper bounds hold in every finite dimension. The lower bound permits a dimension depending on T; Corollary sample-complexity imports Liu, extended-version Theorem 11. No fixed-dimensional lower rate is asserted.



\textbf{Constants.} Constants and sufficient large-budget thresholds may depend on fixed p; they are not uniform as p decreases to 1.



\textbf{Logarithmic factors.} A single logarithm in pointwise sample complexity; its displayed power in risk.



\textbf{Source location.} \href{https://github.com/WU-Wendao/Colombo/blob/main/colombo_sensorium/Complexity_altas/data/sources/colombo/N05/manuscript.pdf}{[1]} Corollary sample-complexity; equation liu-lower; cited Liu Theorem 11; verified labels: eq:liu-lower (line 672), cor:sample-complexity (line 681), eq:sample-complexity (line 683)



Known-index stochastic first-order lower bound, valid for the smaller class of tuning-free families. Every randomized T-query method, even when p is known.



The source lower construction obeys the centered moment convention and allows growing dimension.



\needspace{4\baselineskip}\subsection{B3 - Dyadic portfolio with independent public allocation choices}


\textbf{Admission.} Included in the bracket



\textbf{Direction.} upper



\textbf{Witness or lower-bound scope.} Randomized witness



\AtlasMath{O_p\!\left(\min_{\tau\in\mathcal T_n}\{\sqrt{L\Delta\tau/n}+\sigma\tau^{-(p-1)/p}\}+\sigma M^{1/p}m^{-(p-1)/p}\right)}


\textbf{Sources.} \href{https://github.com/WU-Wendao/Colombo/blob/main/colombo_sensorium/Complexity_altas/data/sources/colombo/N05/manuscript.pdf}{[1]}



\textbf{Conditions.} \AtlasInlineMath{L,\Delta>0}, \AtlasInlineMath{\sigma\ge0}, fixed \AtlasInlineMath{1<p\le2}, integers \AtlasInlineMath{n\ge8,m\ge1}, and \AtlasInlineMath{M(n+m)\le T}.; This nonasymptotic bound remains valid when the analysis-optimal batch size lies outside the interior dyadic range.



\textbf{Dimension conditions.} The upper bound holds for every finite d.



\textbf{Constants.} Constants depend only on fixed p.



\textbf{Logarithmic factors.} No logarithms are suppressed in the displayed allocation formula.



\textbf{Source location.} \href{https://github.com/WU-Wendao/Colombo/blob/main/colombo_sensorium/Complexity_altas/data/sources/colombo/N05/manuscript.pdf}{[1]} Theorem parameter-oracle; equation parameter-oracle; verified labels: thm:parameter-oracle (line 634), eq:parameter-oracle (line 645)



Dyadic portfolio with independent public allocation choices. One common family using L and Delta but not p or sigma.



This nonasymptotic bound remains valid when the analysis-optimal batch size lies outside the interior dyadic range.



\needspace{7\baselineskip}\section{Maintenance appendix: Native proof extensions}


Read-only presentation of the stored native record. Integrity hashes identify source values; this appendix grants no new theorem, source-read credit, review or approval.



Native pointer: /proofs; canonical SHA-256: 4f53cda18c2baa0c0354bb5f\allowbreak{}9a3ecbe5ed12ab4d8e11ba87\allowbreak{}3c2f11161202b945



Native value (empty\_array): []



\needspace{7\baselineskip}\section{References}


[1] Colombo\_N05. Near-Optimal Smooth Nonconvex Optimization with an Unknown Tail Index. 2026. N05\_unknown\_tail\_index\_publication.tex.



\url{https://github.com/WU-Wendao/Colombo/blob/main/colombo_sensorium/Complexity_altas/data/sources/colombo/N05/manuscript.pdf}



[2] Zijian Liu. Online Convex Optimization with Heavy Tails: Old Algorithms, New Regrets, and Applications. 2026. Extended version, March 19, 2026; Theorem 11 as cited by the supplied manuscript.



\url{https://arxiv.org/abs/2508.07473v3}

\end{document}
